\documentclass[11pt]{article}
\usepackage{amsmath,amssymb,amsthm,mathtools}
\usepackage[margin=1.25in]{geometry}

\newtheorem{theorem}{Theorem}[section]
\newtheorem{lemma}[theorem]{Lemma}
\newtheorem{corollary}[theorem]{Corollary}
\newtheorem{proposition}[theorem]{Proposition}
\theoremstyle{definition}
\newtheorem{definition}[theorem]{Definition}
\newtheorem{remark}[theorem]{Remark}

\title{The Gram Point Hierarchy and the Lambert $W$ Generator}
\author{Stephen Crowley}
\date{April 2026}

\begin{document}
\maketitle

\begin{abstract}
The nontrivial zeros of the Riemann zeta function and their relationship
to the Gram points are governed by a hierarchy of phase offset functions
$\psi_k(t) = \vartheta(t) - kt$, where $\vartheta$ is the Riemann--Siegel
theta function. We derive explicit closed-form expressions for the
level-$k$ Gram sequence $\{g_n^{(k)}\}$ in terms of the Lambert $W$
function and show that all correction terms are generated by iterating
a single base inversion operator $\Psi_0$ composed with the Stirling
residual tail of $\vartheta$. The full Lagrange--B\"urmann inversion
produces an exact power series for $g_n^{(k)}$ in which every
coefficient is expressible through
$W\!\bigl((n+\tfrac{1}{8})\,e^{-(2k+1)}\bigr)$, generated by a
recurrence on polynomials $p_j$. The $k$-th level process
$X_k(n) = n\pi - g_n^{(k-1)}$ is thereby expressed entirely in
closed form through the same Lambert $W$ generator at shifted level.
\end{abstract}

\section{Definitions and Setup}

\begin{definition}
The Riemann--Siegel theta function is
\begin{equation}
  \vartheta(t) = \arg\!\left(\pi^{-it/2}\,\Gamma\!\left(\tfrac{1}{4}
    + \tfrac{it}{2}\right)\right),
\end{equation}
normalized so that $\vartheta(t)$ is real and continuous for $t > 0$
with $\vartheta(0) = 0$.
\end{definition}

\begin{definition}
For each integer $k \ge 0$, the $k$-th phase offset function is
\begin{equation}
  \psi_k(t) = \vartheta(t) - kt.
\end{equation}
\end{definition}

\begin{definition}
The Stirling residual tail of $\vartheta$ is the function
$\varphi\colon (0,\infty) \to \mathbb{R}$ defined by the exact identity
\begin{equation}
  \varphi(t) = \operatorname{Im}\!\left[\log\Gamma\!\left(\tfrac{1}{4}
    + \tfrac{it}{2}\right)\right]
    - \frac{t}{2}\log\frac{t}{2} + \frac{t}{2} + \frac{\pi}{8},
  \label{eq:phi-exact}
\end{equation}
so that
\begin{equation}
  \vartheta(t) = \frac{t}{2}\log\frac{t}{2\pi} - \frac{t}{2}
    - \frac{\pi}{8} + \varphi(t).
  \label{eq:theta-exact}
\end{equation}
To verify: from $\vartheta(t) = \operatorname{Im}[\log\Gamma(\tfrac{1}{4}+\tfrac{it}{2})] - \tfrac{t}{2}\log\pi$, substituting~\eqref{eq:phi-exact} into~\eqref{eq:theta-exact} gives
\begin{equation*}
  \frac{t}{2}\log\frac{t}{2\pi} - \frac{t}{2} - \frac{\pi}{8}
  + \operatorname{Im}\!\left[\log\Gamma\!\left(\tfrac{1}{4}+\tfrac{it}{2}\right)\right]
  - \frac{t}{2}\log\frac{t}{2} + \frac{t}{2} + \frac{\pi}{8}
  = \operatorname{Im}\!\left[\log\Gamma\!\left(\tfrac{1}{4}+\tfrac{it}{2}\right)\right]
  - \frac{t}{2}\log\pi,
\end{equation*}
since $\log\frac{t}{2\pi} - \log\frac{t}{2} = -\log\pi$.

The function $\varphi$ is smooth on $(0,\infty)$ and has the asymptotic
expansion
\begin{equation}
  \varphi(t) \sim \frac{1}{48t} + \frac{7}{5760t^3}
    + \frac{31}{80640t^5} + \frac{127}{430080t^7} + \cdots
  \label{eq:phi-asymp}
\end{equation}
arising from the combined contributions of the Stirling remainder
series for $\log\Gamma$ at $z = \tfrac{1}{4}+\tfrac{it}{2}$ and the
subleading terms in the expansion of
$\operatorname{Im}[(z-\tfrac{1}{2})\log z]$ around $z = it/2$.
\end{definition}

\begin{definition}
For each $k \ge 0$, the level-$k$ Gram sequence $\{g_n^{(k)}\}_{n\ge 0}$
is the sequence of positive real numbers satisfying
\begin{equation}
  \psi_k\!\left(g_n^{(k)}\right) = n\pi,
  \qquad n = 0,1,2,\ldots
  \label{eq:gram-def}
\end{equation}
For $k=0$ this gives the ordinary Gram points $g_n^{(0)} = g_n$
satisfying $\vartheta(g_n) = n\pi$.
\end{definition}

\begin{definition}
The $k$-th level process is
\begin{equation}
  X_k(n) = \psi_k\!\left(g_n^{(k-1)}\right)
          = \vartheta\!\left(g_n^{(k-1)}\right) - k\,g_n^{(k-1)},
  \qquad k \ge 1.
  \label{eq:Xk-def}
\end{equation}
\end{definition}

\begin{definition}
The base inversion function is
\begin{equation}
  \Psi_0(c, k) = \frac{2c}{W\!\left(\dfrac{c}{\pi}\,e^{-(2k+1)}\right)},
  \label{eq:Psi0}
\end{equation}
where $W$ is the principal branch of the Lambert $W$ function.
\end{definition}

\section{The Explicit Gram Sequence}

\begin{theorem}[Explicit form of $X_k$]
For every $k \ge 1$ and every $n \ge 0$:
\begin{equation}
  X_k(n) = n\pi - g_n^{(k-1)}.
  \label{eq:Xk-explicit}
\end{equation}
\end{theorem}

\begin{proof}
From definition~\eqref{eq:gram-def} at level $k-1$:
\begin{equation}
  \vartheta\!\left(g_n^{(k-1)}\right) - (k-1)\,g_n^{(k-1)} = n\pi,
\end{equation}
so $\vartheta(g_n^{(k-1)}) = n\pi + (k-1)\,g_n^{(k-1)}$.
Substituting into~\eqref{eq:Xk-def}:
\begin{equation}
  X_k(n) = n\pi + (k-1)\,g_n^{(k-1)} - k\,g_n^{(k-1)} = n\pi - g_n^{(k-1)}.
\end{equation}
\end{proof}

\begin{theorem}[Leading Lambert $W$ solution]
\label{thm:leading}
Set $c = \pi(n + \tfrac{1}{8})$ and $w_k = W\!\bigl(\frac{c}{\pi}\,e^{-(2k+1)}\bigr)$.
The leading-order solution to~\eqref{eq:gram-def} is
\begin{equation}
  g_n^{(k)} = \Psi_0(c, k) + O\!\left(n^{-1}\right)
             = \frac{2\pi\!\left(n + \tfrac{1}{8}\right)}{w_k}
               + O\!\left(n^{-1}\right).
  \label{eq:leading}
\end{equation}
\end{theorem}

\begin{proof}
Substituting~\eqref{eq:theta-exact} into~\eqref{eq:gram-def}:
\begin{equation}
  \frac{t}{2}\log\frac{t}{2\pi} - \frac{(2k+1)}{2}\,t
    = n\pi + \frac{\pi}{8} - \varphi(t).
  \label{eq:implicit}
\end{equation}
Drop $\varphi(t)$ at leading order and set $t = 2\pi e^L$, so that
$\log(t/2\pi) = L$:
\begin{equation}
  \pi e^L\bigl(L - (2k+1)\bigr) = c.
  \label{eq:leading-eq}
\end{equation}
Set $w = L - (2k+1)$, so $L = w + 2k+1$:
\begin{equation}
  \pi e^{2k+1}\,w\,e^w = c,
  \qquad
  w\,e^w = \frac{c}{\pi}\,e^{-(2k+1)}.
\end{equation}
Therefore $w = W\!\bigl(\frac{c}{\pi}\,e^{-(2k+1)}\bigr) = w_k$, giving
$L = w_k + 2k+1$ and
\begin{equation}
  t = 2\pi\,e^{w_k + 2k+1} = 2\pi\,e^{2k+1}\,e^{w_k}.
\end{equation}
Applying the identity $e^{W(z)} = z/W(z)$:
\begin{equation}
  t = 2\pi\,e^{2k+1}\cdot\frac{c/(\pi\,e^{2k+1})}{w_k}
    = \frac{2c}{w_k}.
\end{equation}
The $O(n^{-1})$ error arises from dropping $\varphi(t) = O(t^{-1})$
in~\eqref{eq:implicit}.
\end{proof}

\section{The Lambert $W$ Generator for All Corrections}

\begin{theorem}[Derivatives of $\Psi_0$]
\label{thm:derivatives}
Let $a = \pi\,e^{2k+1}$ and $w = W\!\bigl(\frac{c}{\pi}\,e^{-(2k+1)}\bigr)$.
The $j$-th partial derivative of $\Psi_0(c,k)$ with respect to $c$ is
\begin{equation}
  \partial_c^{\,j}\,\Psi_0(c,k)
    = \frac{2\,e^{-(j-1)w}}{a^{\,j-1}}
      \cdot\frac{p_j(w)}{(1+w)^{2j-1}},
  \label{eq:Psi-deriv}
\end{equation}
where the polynomials $p_j$ satisfy the recurrence
\begin{equation}
  p_1(w) = 1,
  \qquad
  p_{j+1}(w) = p_j'(w)\,(1+w) - \bigl((j-1)w + 3j - 2\bigr)\,p_j(w),
  \label{eq:pj-recurrence}
\end{equation}
with initial values
\begin{equation}
  p_1 = 1,\quad
  p_2 = -1,\quad
  p_3 = w + 4,\quad
  p_4 = -(2w^2 + 14w + 27).
\end{equation}
\end{theorem}

\begin{proof}
From~\eqref{eq:Psi0}, with $z = \frac{c}{\pi}\,e^{-(2k+1)} = c/a$
so that $w = W(z)$ and $c = a\,w\,e^w$:
\begin{equation}
  \partial_c w = \frac{W'(z)}{a} = \frac{1}{a\,e^w(1+w)}.
  \label{eq:dWdc}
\end{equation}
For $j=1$, differentiating $\Psi_0 = 2c/w$:
\begin{equation}
  \partial_c\,\Psi_0
    = \frac{2}{w} - \frac{2c}{w^2}\cdot\frac{1}{a\,e^w(1+w)}
    = \frac{2}{w} - \frac{2\,a\,w\,e^w}{w^2\,a\,e^w(1+w)}
    = \frac{2}{w} - \frac{2}{w(1+w)}
    = \frac{2}{1+w},
\end{equation}
confirming~\eqref{eq:Psi-deriv} with $p_1 = 1$ and the convention
$e^{0\cdot w} = 1$.

Now assume~\eqref{eq:Psi-deriv} holds at index $j$.
Differentiating
$\partial_c^{\,j}\,\Psi_0 = 2\,a^{1-j}\,e^{-(j-1)w}\,p_j(w)\,(1+w)^{-(2j-1)}$
with respect to $c$ via the product rule, with each $w$-dependent
factor differentiated using~\eqref{eq:dWdc}:
\begin{align}
  \partial_c^{\,j+1}\,\Psi_0
  &= \frac{2}{a^{j-1}}\,\frac{e^{-(j-1)w}}{(1+w)^{2j-1}}
     \Bigl[p_j'(w) - (j-1)\,p_j(w)
           - \frac{(2j-1)\,p_j(w)}{1+w}\Bigr]
     \cdot\frac{1}{a\,e^w(1+w)} \notag\\[4pt]
  &= \frac{2\,e^{-jw}}{a^{\,j}}\cdot
     \frac{p_j'(w)(1+w) - \bigl((j-1)(1+w) + (2j-1)\bigr)\,p_j(w)}
          {(1+w)^{2j+1}}.
  \label{eq:induction-step}
\end{align}
Since $(j-1)(1+w) + (2j-1) = (j-1)w + 3j - 2$, defining
$p_{j+1}(w) = p_j'(w)(1+w) - ((j-1)w + 3j-2)\,p_j(w)$
gives~\eqref{eq:pj-recurrence} and establishes~\eqref{eq:Psi-deriv}
by induction.

Explicitly:
\begin{align}
  p_2 &= 0\cdot(1+w) - 1\cdot 1 = -1, \\
  p_3 &= 0\cdot(1+w) - (w+4)\cdot(-1) = w+4, \\
  p_4 &= (1+w) - (2w+7)(w+4) = 1+w - 2w^2-15w-28 = -(2w^2+14w+27).
\end{align}
\end{proof}

\begin{theorem}[Exact Gram sequence via Lagrange--B\"urmann inversion]
\label{thm:lagrange}
Define the correction function
\begin{equation}
  \eta_k(c) = \varphi\!\left(\Psi_0(c,k)\right),
\end{equation}
where $\varphi$ is defined by~\eqref{eq:phi-exact}. Then the exact
solution to~\eqref{eq:gram-def} is
\begin{equation}
  g_n^{(k)} = \sum_{j=0}^{\infty} \frac{(-1)^j}{j!}\,
    \frac{d^{\,j-1}}{dc^{\,j-1}}\!\left[
      \partial_c\,\Psi_0(c,k)\cdot\eta_k(c)^j
    \right]_{c\,=\,\pi(n+1/8)},
  \label{eq:lagrange}
\end{equation}
where $\frac{d^{-1}}{dc^{-1}}$ at $j=0$ denotes the identity
operator, and all derivatives of $\Psi_0$ are given by~\eqref{eq:Psi-deriv}.
\end{theorem}

\begin{proof}
Equation~\eqref{eq:gram-def} with~\eqref{eq:theta-exact} is
\begin{equation}
  f(t) \coloneqq \frac{t}{2}\log\frac{t}{2\pi}
    - \frac{(2k+1)}{2}\,t = c - \varphi(t),
  \label{eq:fp}
\end{equation}
with $c = \pi(n + \tfrac{1}{8})$.
The function $\Psi_0(\,\cdot\,,k)$ is the exact inverse of $f$:
$f(\Psi_0(c,k)) = c$ for all $c$ in the range of $f$.
Equation~\eqref{eq:fp} is therefore the fixed-point equation
\begin{equation}
  t = \Psi_0\!\bigl(c - \varphi(t),\,k\bigr).
  \label{eq:fp2}
\end{equation}
By the Lagrange--B\"urmann theorem (in the form given by
Whittaker--Watson, \S7.32): if $t = \Psi_0(c - \varphi(t),k)$
and $t_0 = \Psi_0(c,k)$, then for any analytic $F$,
\begin{equation}
  F(t) = F(t_0) + \sum_{n=1}^{\infty}\frac{(-1)^n}{n!}\,
    \frac{d^{\,n-1}}{dc^{\,n-1}}\!\left[
      F'(t_0)\,\Psi_0'(c,k)\,\eta_k(c)^n
    \right],
\end{equation}
where $\eta_k(c) = \varphi(\Psi_0(c,k))$. Taking $F$ to be the
identity gives~\eqref{eq:lagrange}.

Convergence holds for $c$ in a disk around $\pi(n+\tfrac{1}{8})$
of radius equal to the distance from $c$ to the nearest singularity
of $\Psi_0(\,\cdot\,,k)$ in the complex $c$-plane. For the principal
branch of $W$, this singularity occurs at $c = -a/e$ (where
$a = \pi\,e^{2k+1}$), giving radius $|c + a/e| > c > 0$ for $n \ge 0$.
Since $|\eta_k(c)| = O(c^{-1})$ while this radius grows linearly
in $c$, the series converges for all $n \ge 0$.
\end{proof}

\begin{corollary}[First correction term and second-order structure]
\label{cor:corrections}
Writing $w = w_k = W\!\bigl((n+\tfrac{1}{8})\,e^{-(2k+1)}\bigr)$,
$\Psi = \frac{2\pi(n+1/8)}{w}$, $a = \pi\,e^{2k+1}$,
and $\eta = \varphi(\Psi)$:
\begin{equation}
  g_n^{(k)} = \Psi
    - \frac{2}{1+w}\cdot\eta
    + \frac{1}{2}\frac{d}{dc}\!\left[\frac{2}{1+w}\cdot\eta^2\right]_{c\,=\,\pi(n+1/8)}
    + O\!\left(\eta^3\right),
  \label{eq:two-terms}
\end{equation}
where the $j = 2$ term expands via the product rule to
\begin{equation}
  \frac{1}{2}\frac{d}{dc}\!\left[\frac{2}{1+w}\cdot\eta^2\right]
  = -\frac{e^{-w}\,\eta^2}{a\,(1+w)^3}
    + \frac{4\,\eta\,\varphi'(\Psi)}{(1+w)^2},
  \label{eq:j2-expanded}
\end{equation}
with $\varphi'(\Psi) = \frac{d}{dt}\varphi(t)\big|_{t=\Psi}$. Both summands
in~\eqref{eq:j2-expanded} are $O(\eta/\Psi)$ and contribute at the same
asymptotic order.
\end{corollary}

\begin{proof}
The $j=0$ and $j=1$ terms of~\eqref{eq:lagrange} give
$\Psi - \frac{2}{1+w}\,\eta$. The $j=2$ term is
$\frac{1}{2}\frac{d}{dc}[\partial_c\Psi_0\cdot\eta^2]$.
Applying the product rule:
\begin{equation}
  \frac{1}{2}\frac{d}{dc}\!\left[\frac{2}{1+w}\cdot\eta^2\right]
  = \frac{1}{2}\cdot\partial_c^2\Psi_0\cdot\eta^2
    + \frac{2}{1+w}\cdot\eta\cdot\eta'(c).
\end{equation}
From Theorem~\ref{thm:derivatives} at $j=2$:
$\partial_c^2\Psi_0 = -\frac{2\,e^{-w}}{a\,(1+w)^3}$.
And $\eta'(c) = \varphi'(\Psi)\cdot\partial_c\Psi_0
= \varphi'(\Psi)\cdot\frac{2}{1+w}$. Substituting:
\begin{equation*}
  \frac{1}{2}\cdot\!\left(\!-\frac{2\,e^{-w}}{a(1+w)^3}\right)\eta^2
  + \frac{2}{1+w}\cdot\eta\cdot\varphi'(\Psi)\cdot\frac{2}{1+w}
  = -\frac{e^{-w}\,\eta^2}{a(1+w)^3} + \frac{4\,\eta\,\varphi'(\Psi)}{(1+w)^2},
\end{equation*}
yielding~\eqref{eq:j2-expanded}.
\end{proof}

\begin{corollary}[Explicit $k$-th level process]
The $k$-th level process $X_k(n) = n\pi - g_n^{(k-1)}$ is
\begin{equation}
  X_k(n) = n\pi
    - \sum_{j=0}^{\infty} \frac{(-1)^j}{j!}\,
      \frac{d^{\,j-1}}{dc^{\,j-1}}\!\left[
        \frac{2}{1+w_{k-1}}
        \cdot \eta_{k-1}(c)^j
      \right]_{c\,=\,\pi(n+1/8)},
  \label{eq:Xk-full}
\end{equation}
with $w_{k-1} = W\!\bigl((n+\tfrac{1}{8})\,e^{-(2k-1)}\bigr)$
and $\eta_{k-1}(c) = \varphi(\Psi_0(c,k-1))$.
For $k=1$:
\begin{equation}
  X_1(n) = n\pi
    - \frac{2\pi\!\left(n+\tfrac{1}{8}\right)}{W\!\left((n+\tfrac{1}{8})/e\right)}
    + \frac{2}{1+w_0}\cdot\varphi\!\left(
        \frac{2\pi(n+\tfrac{1}{8})}{w_0}\right)
    + O\!\left(\varphi^2\right),
  \label{eq:X1-explicit}
\end{equation}
with $w_0 = W\!\bigl((n+\tfrac{1}{8})/e\bigr)$.
\end{corollary}

\begin{proof}
Substitute~\eqref{eq:lagrange} with $k$ replaced by $k-1$
into~\eqref{eq:Xk-explicit}.
\end{proof}

\begin{remark}
The polynomials $p_j$ in~\eqref{eq:pj-recurrence} depend only on $w$
and not on $k$. However, the factor $e^{-(j-1)w}$
in~\eqref{eq:Psi-deriv} introduces additional $w$-dependence beyond
the polynomial, so the full $j$-th derivative is not purely a rational
function of $w$ and $(1+w)$. The dependence on the level $k$ enters
through the argument of $W$: replacing $k$ by $k+1$ multiplies the
argument by $e^{-2}$. Every level of the hierarchy is therefore
generated by the same algebraic structure --- the
recurrence~\eqref{eq:pj-recurrence} and the base function
$\Psi_0$ --- evaluated at a geometrically rescaled argument of $W$.
\end{remark}

\end{document}
